\documentclass[10pt,a4paper]{amsart}
\usepackage[margin=19mm]{geometry}
\usepackage{amsmath,amssymb,mathtools}
\usepackage[scr=rsfs,cal=euler]{mathalfa}
\usepackage{xfrac}
\usepackage[unicode,hidelinks,bookmarks=false]{hyperref}
\newcommand{\ie}{i.e.\ }
\newcommand{\cf}{cf.\ }
\newcommand{\resp}{respectively\ }
\newcommand{\Subsection}{\subsection}
\providecommand{\nobreakdash}{\nobreak\hbox{-}}
\setlength{\emergencystretch}{2em}
\multlinegap0pt
\usepackage{bm}
\usepackage{bbm}
\usepackage{dsfont}
\usepackage{xspace}
\usepackage{cancel}
\usepackage{mathtools}
\usepackage{amsthm}
\makeatletter
\@ifundefined{theorem}{\newtheorem{theorem}{Theorem}}{}
\@ifundefined{corollary}{\newtheorem{corollary}[theorem]{Corollary}}{}
\@ifundefined{proposition}{\newtheorem{proposition}[theorem]{Proposition}}{}
\makeatother
\title[Hopf-Galois structures on separable field extensions of degr]{Hopf-Galois structures on separable field extensions of degree related to Cunningham chains}
\author{Andrew Darlington}
\date{Source: arXiv:2508.03384v1 (CC BY 4.0)}
\begin{document}
\maketitle

\noindent\emph{Source and licence.} Hopf-Galois structures on separable field extensions of degree related to Cunningham chains. Version arXiv:2508.03384v1 (CC BY 4.0), distributed under the Creative Commons Attribution 4.0 International licence, as recorded in the arXiv licence metadata for this exact version and shown on the abstract page https://arxiv.org/abs/2508.03384v1. Full rights notice: licenses/arxiv-2508.03384-CC-BY-4.0.txt.\par\smallskip

\noindent\textit{Editorial scope.} Counted unit: the Proposition whose source label is \texttt{2-subgroups} in the article (source0.tex lines 200--203, source label \texttt{2-subgroups}), delivered together with the article's own complete original proof of it (source0.tex lines 204--262, one \texttt{proof} environment of 59 lines). No mathematics is written, repaired or re-derived: statement and proof are the article's own text line for line. Required context retained uncounted: c\_2\_subgroups (lines 154--157). Every definition, display and label of those lines is retained unchanged. Omitted from this package and not redistributed: 1--153, 158--199, 263--1023. Nothing inside a retained excerpt is omitted or altered: every retained body is delivered exactly as the article's lines are written, so no omission receipt of any kind accompanies this package. Citations: the retained excerpts carry no citation command at all, so no bibliography entry of the article is transcribed and no reference is redistributed. Reference adaptation: none was needed and none was made; every reference inside the delivered unit resolves inside it, so no unresolved reference or citation is produced. Printed slips kept literal: no printed word, symbol, hypothesis or number of the counted statement or of its proof was altered. Layout and class: the article's own class is \texttt{amsart} with packages including inputenc, amsmath, amsfonts, amssymb, pifont, enumerate, enumitem, wasysym, mathrsfs, xfrac, array, siunitx, mathtools, amsthm; the wrapper is the lane's pinned \texttt{amsart} editorial wrapper at 10pt on A4, which restates the theorem-like environments the retained text needs and adds the article's own macro definitions that the retained text needs (none), with extra wrapper packages (bm, bbm, dsfont, xspace, cancel, mathtools). The delivered numbering is the unit's own. Assets: no retained span contains a figure environment, a graphics-inclusion command, a PDF-page inclusion, a table or a TikZ picture; no image file is copied into this package, the package declares no asset dependency and no image file is redistributed with it. Licence: version arXiv:2508.03384v1 of this article is distributed under the Creative Commons Attribution 4.0 International licence, as recorded in the arXiv licence metadata for this exact version and shown on the abstract page https://arxiv.org/abs/2508.03384v1. A full rights notice is delivered at licenses/arxiv-2508.03384-CC-BY-4.0.txt. Characters: every retained excerpt is pure ASCII, so the delivered engine, XeLaTeX with its default Latin Modern text font, reports no missing character.
\begin{corollary}\label{c_2_subgroups}
    The number of subgroups of $(C_2)^l$ is
    \[\sum_{k=0}^l\prod_{i=1}^k\frac{2^{l-(i-1)}-1}{2^i-1}.\]
\end{corollary}

\begin{proposition}\label{2-subgroups}
The number of subgroups of $\langle \beta_1,\ldots,\beta_{l-1},\gamma\rangle$ is
\[\Sigma_l:=\sum_{k=0}^{l-1}(2^{l-(k+1)}x+1)\prod_{i=1}^k\frac{2^{l-i}-1}{2^i-1}.\]
\end{proposition}

\begin{proof}
    Let $B:=\langle \beta_1,\ldots,\beta_{l-1},\gamma\rangle$ and $B':=\langle \beta_1,\ldots,\beta_{l-1}\rangle$. We first determine the form of an arbitrary subgroup $\langle z_1,\ldots, z_r \rangle$ of $B$. Suppose then that $z_i=\beta_1^{a_{i1}}\beta_2^{a_{i2}}\cdots\beta_{l-1}^{a_{il-1}}\gamma^{a_{il}}$. We can consider the following matrix
    \[\begin{pmatrix}
        a_{11} & a_{12} & \cdots & a_{1l}\\
        a_{21} & a_{22} & \cdots & a_{2l}\\
        \vdots & \vdots & \ddots & \vdots\\
        a_{r1} & a_{r2} & \cdots & a_{rl}
    \end{pmatrix}\]
    where each row represents a generator of the subgroup. The entries can be seen to be in $\mathbb{Z}/2^x\mathbb{Z}$, but we have $a_{ij} \in \{0,1\}$ for $j<l$. We note that applying elementary row operations does not change the subgroup, and so we may assume without loss of generality that $a_{ii}=\gcd(a_{ii},a_{(i+1)i},\ldots,a_{ri})$ and then that $a_{ij}=0$ for each $1 \leq i \leq l$ and $i<j \leq r$. We obtain the following $l \times l$ matrix, which is in row echelon form:
    \[\begin{pmatrix}
        a_{11} & a_{12} & \cdots & a_{1l}\\
        0 & a_{22} & \cdots & a_{2l}\\
        \vdots & \vdots & \ddots & \vdots\\
        0 & 0 & \cdots & a_{ll}
    \end{pmatrix}\]
    where we delete rows $l+1,\ldots,r$, which are all zero. Thus an arbitrary subgroup of $B$ has the form
\begin{equation}\label{subgroup}
    \langle \beta'_1\gamma^{2^{x-c_1}},\ldots,\beta'_{l-1}\gamma^{2^{x-c_{l-1}}},\gamma^{2^{x-c_l}}\rangle
\end{equation}
    for some fixed $0 \leq c_i \leq x$, $1\leq i \leq l$, and some $\beta'_j \in B'$, $1 \leq j \leq l-1$. We note further that if we have two generators of the subgroup $g_1:=\beta'_i\gamma^{2^{x-c_i}}, g_2:=\beta'_j\gamma^{2^{x-c_j}}$ such that $0 \leq c_i \leq c_j$, then $g_2^{-2^{c_j-c_i}}g_1=\beta'_i$, and so we may assume that $c_i=0$. More generally therefore, we may assume that $c_i=c_j$ for all nonzero $c_i,c_j$ in (\ref{subgroup}). Thus we now denote $c_i$ simply as $c$.
 
    We now count the subgroups $S$ of this form for a fixed $c$ with $1 \leq c \leq x$.

    (i$_c$) If $\gamma^{2^{x-c}}$ is a generator of $S$, then we can assume all the other generators involve only $\beta_1,\ldots, \beta_{l-1}$. So $S$ corresponds to a matrix of $l$ columns in reduced row echelon form (RREF) such that the first $l-1$ columns have entries in $\{0,1\}$ and the final column has all zeros apart from $2^{x-c}$ in the last entry.
   
    (ii$_c$) If $\gamma^{2^{x-c}}$ does not occur as a generator of $S$, then $S$ corresponds to a matrix of $l$ columns in RREF such that the first $l-1$ columns have entries in $\{0,1\}$ and entries in the final column are either $0$ or $2^{x-c}$, with $0$ in the last entry.
    
    We note that the cases (i$_c$) are disjoint from (i$_{c'}$) for $c \neq c'$ and from (ii$_{c''}$) for any $1 \leq c'' \leq x$. Further, we note that the matrix determines the value of $c$ except when the last column consists entirely of zeros; in this case the matrix corresponds to a subgroup of $B'$. Thus, given two values of $c$, say $c$ and $c'$, the cases (ii$_c$) and (ii$_{c'}$) are not disjoint, but intersect precisely when every entry of the final column of the matrix is equal to 0.
    
    Now, for each $c>0$, we have a bijection between the matrices in (i$_c$) and the matrices in (ii$_c$), and the matrices with $l$ columns over $\mathbb{F}_2$ in RREF by replacing each occurrence of $2^{x-c}$ by $1$, whenever it occurs. The number of these matrices is the same as the number of subgroups of $(C_2)^l$. That is, by Corollary \ref{c_2_subgroups}:
    \begin{equation}\label{l-RREF}
        \sum_{k=0}^l\prod_{i=1}^k\frac{2^{l+1-1}-1}{2^i-1}.
    \end{equation}
    To obtain the total number of subgroups of $B$, we note that we can just sum (\ref{l-RREF}) over $1 \leq c \leq x$, with the caveat that we over-count by
    \[(x-1)\sum_{k=1}^{l-1}\prod_{i=1}^k\frac{2^{l-i}-1}{2^i-1}.\]
    This is precisely due to the intersection between the two cases (ii$_c$) and (ii$_{c'}$) described above. In theory, the subgroups of $B'$ correspond to the case $c=0$, but this is already covered as discussed. So we need to take away $(x-1)$ times the number of subgroups of $B' \cong (C_2)^{l-1}$ from our initial count, as we need to count it exactly once. By Corollary \ref{c_2_subgroups}, the number of subgroups of $B'$ is precisely
    \[\sum_{k=1}^{l-1}\prod_{i=1}^k\frac{2^{l-i}-1}{2^i-1}.\]
    Therefore, the number of subgroups of $B$ is:
    \begin{equation}\label{subsofB}
        x\sum_{k=0}^l\prod_{i=1}^k\frac{2^{l+1-i}-1}{2^i-1}-(x-1)\sum_{k=0}^{l-1}\prod_{i=1}^k\frac{2^{l-i}-1}{2^i-1}.
    \end{equation}
    We claim that this simplifies to
    \[\sum_{k=0}^{l-1}(2^{l-(k+1)}x+1)\prod_{i=1}^k\frac{2^{l-i}-1}{2^i-1}.\]
    To see this, we may first split the sum (\ref{subsofB}) as follows:
    \begin{equation}\label{split_sum}
        x\left[\sum_{k=0}^l\prod_{i=1}^k\frac{2^{l+1-i}-1}{2^i-1}-\sum_{k=0}^{l-1}\prod_{i=1}^k\frac{2^{l-i}-1}{2^i-1}\right]+\sum_{k=0}^{l-1}\prod_{i=1}^k\frac{2^{l-i}-1}{2^i-1}.
    \end{equation}
    Noting that the product
    \[\prod_{i=1}^k\frac{2^{l+1-i}-1}{2^i-1}\]
    evaluates to $1$ when $k=l$, we have that the `$x$' part of (\ref{split_sum}) can be written as
    \[1+\sum_{k=0}^{l-1}\left[\prod_{i=1}^k\frac{2^{l+1-i}-1}{2^i-1}-\prod_{i=1}^k\frac{2^{l-i}-1}{2^i-1}\right].\]
    We further note that the summand for $k=0$ evaluates to $1-1=0$, and therefore we may start the sum from $k=1$. Considering, then, the difference between the two products in the summand, this evaluates to
    \[1+\sum_{k=1}^{l-1}\frac{2^l-2^{l-k}}{2^k-1}\prod_{i=1}^{k-1}\frac{2^{l-i}-1}{2^i-1}.\]
    Noting now that $(2^l-2^{l-k})/(2^k-1)=2^{l-k}$ and making the change of variables $k \mapsto k+1$, we see that this is equal to
    \[1+\sum_{k=0}^{l-2}2^{l-(k+1)}\prod_{i=1}^k\frac{2^{l-i}-1}{2^i-1}.\]
    Finally, noting that when $k=l-1$, the summand evaluates to $1$, and hence the `$x$' part of the sum can be written as
    \[\sum_{k=1}^{l-1}2^{l-(k+1)}\prod_{i=1}^k\frac{2^{l-i}-1}{2^i-1}.\]
    The rest is clear.
\end{proof}

\end{document}
